\section{Introduction}
\label{sec:introduction}

\PARstart{A}{n} unmeasured load changes. Several PMUs respond, a detector raises an alarm, and the network keeps oscillating after the load is restored. A second intervention then begins during that recovery. The operational question is no longer whether one signal exceeds a threshold: which physical intervention explains the reports observed so far, where did it originate, and when is that explanation distinguishable from the alternatives?

\begin{figure*}[t]\centering
\resizebox{\textwidth}{!}{\begin{tikzpicture}[font=\sffamily\footnotesize\mdseries,>=Latex,
 box/.style={draw=black!40,rounded corners=3pt,align=left,text width=4.45cm,minimum height=1.48cm,inner sep=6pt},
 arr/.style={->,thick,draw=black!65}]
\node[box,fill=accessblue!8] (y) at (0,0) {\textbf{1. Receive a PMU report}\\Eight PMUs: $V$, terminal $I$, $f$, ROCOF.\\Observed history through time $t$.};
\node[box,fill=green!7] (physics) at (5.35,0) {\textbf{2. Predict a physical response}\\DAE dynamics and network constraints.\\Oriented terminal-current map.};
\node[box,fill=orange!9] (hyp) at (10.7,0) {\textbf{3. Compare physical histories}\\New load/line or existing recovery.\\Replay from the estimated past state.};
\node[box,fill=orange!9] (like) at (10.7,-2.7) {\textbf{4. Compare prediction and data}\\Filter state and colored noise.\\Integrate severity; score a report once.};
\node[box,fill=green!7] (mass) at (5.35,-2.7) {\textbf{5. Retain active and cause masses}\\Mix fixed network models into $M_c(t)$.\\$\Delta M$: evidence change + reweighting.};
\node[box,fill=accessblue!8] (out) at (0,-2.7) {\textbf{6. Emit a causal decision}\\Alarm/diagnosis emission times.\\Source candidates and restoration.};
\draw[arr] (y)--(physics);\draw[arr] (physics)--(hyp);\draw[arr] (hyp)--(like);\draw[arr] (like)--(mass);\draw[arr] (mass)--(out);
\node[draw=black!35,fill=black!3,rounded corners=2pt,align=center,text width=15.1cm,minimum height=.68cm] (eval) at (5.35,-4.8) {\textbf{Evaluation only:} seal predictions, then compare with reserved event truth. Physical OFF and electrical recovery are different times.};
\draw[arr,dashed] (out.south)--++(0,-.45)-|(eval.west);
\end{tikzpicture}
}
\caption{The complete scientific argument. Physics maps an intervention at an unmeasured source into predicted PMU trajectories; causal evidence chooses among those explanations. Event truth enters only evaluation. The figure describes the implemented restricted load/line component, not an already validated all-family joint estimator.}
\label{fig:joint-architecture}\end{figure*}

Sparse PMUs make the question difficult because a source need not have a sensor at its own bus. Its voltage, current and frequency effects propagate through the network \cite{phadke2008synchronized,zhao2019dse}. Uncertain operating context and network parameters can also change those effects. Moreover, physical OFF is not electrical recovery. A method that resets at restoration can treat the remaining response as another event or lose a real event that begins before the oscillation decays.

Strong existing methods solve parts of this problem. Dynamic state estimators combine power-system equations and PMU histories \cite{katanic2024recursive}; supervised classifiers learn sparse-PMU event/source labels \cite{pandey2020realtime,yildiz2024few}; physical outage detectors compare network hypotheses sequentially \cite{chen2016quickest,rovatsos2017outage}; and high-fidelity fault trajectories already support sparse-PMU diagnosis \cite{warner2023highfidelity}. The question here is how to retain an explanation through restoration and compare its recovery with a later intervention. A low state error does not prove source identity, and a score on represented assets does not prove source transfer.

Our approach uses the physical intervention as the object of inference.
DAE derivatives construct switched predictors for load and line candidates.
Starting from an estimated past state, those predictors describe what the
eight PMUs should report. Conditional likelihoods compare the predictions
with observations while accounting for severity and instrumental memory.
A finite network-model bank represents parameter uncertainty. Active and
originating-cause masses distinguish an ongoing intervention from the
response that remains after restoration.

The paper makes three scoped contributions:
\begin{enumerate}
\item \textbf{Recovery-aware physical attribution:} an executable causal load/line history bank starts from estimated past states, retains the originating cause after physical OFF, and compares recovery with new-intervention histories under fixed network hypotheses. We specify its evidence and publication rules.
\item \textbf{A qualified source-resolution analysis:} a separate archived load-support study exposes ambiguity and the different information used by joint and conditional source decisions. A prior-consistency audit distinguishes implementation effects from a claimed physical limit; this study is not counted as causal accuracy.
\item \textbf{A reproducible operational test:} the nominal/mixture comparison and a failed admission ablation expose the trade-off between timely source attribution and extra alarms. Separate state, batch and stream contracts retain their own denominators. R11 and nominal remain frozen co-references for confirmation.
\end{enumerate}

The central question is whether this recovery-aware construction preserves
timely source decisions without excessive extra alarms. The experiments
separate development from independent evaluation and keep detection,
eventual attribution and attribution before restoration as different
outcomes. Physical templates, onset scanning and Bayesian updates have
close precedents; the claimed contribution is their specified operational
construction for the declared load/line problem and its measured limits.


